% The integrated commentary of the concise study, beginning on physical page
% 5. It is organized by the questions the formulary raises across its
% elements, not by the order of the Mass, and it sets the three readings'
% witnesses beside one another where they answer the same question. Every
% claim is the expansive study's; the fuller argument, its quotations and its
% apparatus remain there.

\runninghead{Detailed Commentary}
\section{The Propers: Detailed Commentary}
\zlabel{triptych:concise:commentary:start}

The formulary can be read whole in three ways. The first, the Lamb who
receives the kingdom, reads the Mass from the Introit and the Offertory. The second, the kingdom not from hence and therefore everywhere,
reads it from the Gospel and the chants of universal dominion. The third, the
kingdom handed over that has no end, reads it from the Preface, the Alleluia
and the Communion. They answer different questions and are complementary. Their witnesses
divide on particular verses: the decree of the second psalm, the birth of
Jn 18:37, the Son's subjection to the Father, and the moment when the Lord
sits King for ever.

\subsection{What the Lamb receives, and why the Son must ask}

The difficulty lies in two verbs. The Introit's slain Lamb is worthy
\latin{accipere \dots\ divinitatem}, to receive divinity; the Offertory's Son,
just declared begotten, is told \latin{Postula a me}, ask of me. If he is the
eternal Son, he can neither receive divinity nor lack the nations. St~Jerome,
expounding the verse of Daniel from which the Alleluia comes, writes that the
Son of man's receiving of power, honour and kingdom must be understood by the
Apostle's ``Who being in the form of God \dots\ emptied himself, taking the
form of a servant \dots\ becoming obedient unto death'' (Phil 2:6--8), which
the Arians, he says, should have read piously (\work{In Dan.}\ II, on
7:13--14). The same key serves every text of receiving: read with it, such a
text speaks of what the eternal Son receives in the nature he took. The
Fathers answer the question where the second psalm and Daniel raise it, and
their answers carry over to the Lamb, who is worthy to receive because he was
slain.

St~Hilary of Poitiers reads the decree of Ps 2:7 of the resurrection, on
apostolic authority: Paul at Antioch had said that God fulfilled the second
psalm's ``Thou art my Son: this day have I begotten thee'' by ``raising up
Jesus'' (Acts 13:33). The ``today'' of the decree is thus the begetting of the
firstborn from the dead (\work{Tractatus super Psalmos} 2.30). The risen Lord's
``All power is given to me'' (Mt 28:18) shows that he asked for what he
received, and Hilary concludes, \latin{accepit ergo hereditatem gentium, quam
poposcit}, he received the inheritance of the nations, which he asked for
(2.31); he finds the asking itself in the prayer of the Supper, ``Glorify thy
Son'' (Jn 17:1). What was asked he answers from Philippians: the one who took
the form of a servant is raised into the name above every name, \latin{in
deum namque, quia nullum ultra deum nomen est, prouehitur, eique hoc petenti,
id est, ut esset, quod fuerat ante, donatur}, raised into God, since there is
no name beyond God, and given, at his asking, to be what he had been before
(2.33). The Son gains nothing as God; the servant he became is raised into the
glory the Word always had. And the inheritance is not territory but the
nations baptized and taught, carried \latin{ex iniusto \dots\ iure
dominantium in regnum aeternum diuinumque}, out of the unjust right of those
who lorded it over them into an eternal and divine kingdom, the inheritance
\latin{primogeniti huius ex mortuis} (2.31). Hilary is expounding the psalm,
but his words are the Epistle's: those transferred ``into the kingdom of the
Son of his love,'' and the Son ``the firstborn from the dead'' (Col 1:13, 18).

St~Augustine divides the same two verses differently. ``To-day have I begotten
Thee'' has ``a divine interpretation'': it is the eternal generation, and
``today'' signifies the presence of eternity (\work{Enarrationes in Psalmos}
2.6). Only ``Ask of Me'' is temporal, spoken of ``the Manhood which He took on
Himself, who offered up Himself as a Sacrifice in the stead of all
sacrifices'' (2.7). Hilary and Augustine therefore agree on the asking, which
belongs to the Son made man, and part on the decree, which Hilary reads of the
resurrection and Augustine of eternity. St~Thomas Aquinas turns Augustine's
division into a principle. Christ holds the nations \latin{duplici jure:
scilicet hereditario, et meritorio}: by inheritance, from the eternal
generation, \latin{hodie genui te, idest aeternaliter}; and by merit, because
God wills that his gifts be received by asking, and \latin{hoc exemplum voluit
ostendere per Christum, quia voluit quod peteret, quod sibi jure hereditario
competebat}, he willed to show this example in Christ, that he should ask for
what was his by right of inheritance. Christ asks by his prayer, ``for them
also who through their word shall believe in me'' (Jn 17:20), and by his
passion; and he holds the nations not as Peter or Paul held them, as
ministers, but as lord (\work{Postilla super Psalmos} 2, nn.~5--6). The Son
does not ask because he lacks; he asks because the Father wills that even
what is owed be received as a gift, and the asking is his passion.
Cassiodorus states the rule at its widest: whenever Christ asks that something
be given him, \latin{naturam humanitatis evidenter ostendi}, the nature of his
humanity is plainly shown, for \latin{natura humanitatis accepit, quod semper
divina possedit} (\work{Expositio Psalmorum}, on Pss 71 and 2). Theodoret of Cyrus says
the same of Ps 2:8, \latin{etiam ut homo universorum dominium accipit}, but
parts from Augustine on the decree, which he will not apply to the Godhead,
since of that generation the Father says ``from the womb before the day star
I begot thee'' (Ps 109:3).

St~Bede brings the same grammar to the Apocalypse itself. The one whom the
elder calls a Lion and John sees as a Lamb is one Lord, a Lamb in dying
innocently and a Lion in conquering death; and when the Lamb takes the book,
\latin{dispensationem ipsam Incarnationis a Patre et a seipso, secundum id quod
Deus est, dispositam accepisse Filius hominis dicitur}, the Son of man is said
to receive the very dispensation of the Incarnation, which the Father and he
himself, as God, had disposed (\work{Explanatio Apocalypsis} I, on 5:5--7). The
giver and the receiver are the same Son in his two natures. Read with
Hilary's servant raised into the name of God and Bede's Son of man who
receives what he disposed as God, the Introit's \latin{accipere \dots\
divinitatem} means, on the first reading, that the slain Lamb receives, in the
nature he took, the name and glory of the God he always was. He does not
become what he was not.

The strongest objection comes from this Mass's Gospel. On ``For this was I
born,'' St~John Chrysostom writes: ``If then He was born a king, all His other
attributes are by Generation, and He has nothing which He received in
addition''; and he draws a rule for every text of giving: when you hear that
the Father ``has given to the Son also to have life,'' ``deem of nothing else
but His generation, and so of the rest'' (\work{Homilies on John} 83.4; Jn
5:26). St~Cyril of Alexandria says of the same verse that the sovereignty is
``not \dots\ added unto it from without'' (\work{Commentary on John} XII).
Augustine's division grants their premise, since the kingship by generation is
the decree and is not received. It does not meet Chrysostom's rule. His rule
refers a text of giving to the eternal generation; the rule Cassiodorus states
and Theodoret applies refers the asking and receiving to the humanity. The
witnesses agree that Christ is King by nature, and on how to read the texts of
giving addressed to the Son they disagree. Hilary's resurrection reading is
compatible with the answer only because he does not make the resurrection the
beginning of the Sonship: what is given is \latin{ut esset, quod fuerat
ante}. \work{Quas primas} holds both titles together, a dominion by essence
and nature and a second acquired by redemption, at the price of his precious
blood (1~Pet 1:18--19; n.~13).

The answer changes how the surrounding elements are heard. The Alleluia keeps
only the possessing half of Dan 7:14, the power as held; the giving that
Jerome expounds by Philippians stands in the half that is not sung. The
Epistle supplies the means and the moment of the receiving, redemption
``through his blood'' and the head of the Church as man, \latin{secundum quod
huiusmodi}, in Aquinas's gloss (\work{Super Col.}\ 1 lect.~5). The Secret then
names the asking that the Fathers found in the psalm. The sacrifice by which,
in Augustine's words, he ``offered up Himself as a Sacrifice in the stead of
all sacrifices'' is offered now, \latin{quem sacrificiis praesentibus
immolamus}, and what is asked of him is what Hilary's inheritance consists of,
the nations' unity and peace.

\subsection{A prisoner whose kingdom is everywhere}

Read from the first question, the Gospel's prisoner is acquiring his dominion
in the trial itself: he is captive, Chrysostom says, ``by consenting, and
giving Himself up'' (\work{Homilies on John} 83.4), the passion already begun.
Read from the second, the same prisoner says where his kingdom comes from, and
the Mass's chants of universal dominion depend on the answer.

St~Augustine begins from Jesus' counter-question, ``Sayest thou this thing of
thyself?'', put so that the vain notion of his kingdom that Gentiles and Jews
alike had formed might be shown up (\work{Tractates on the Gospel of John}
115.1). Then he gives the answer its widest audience: ``hear, all ye kingdoms
of the earth: I interfere not with your government in this world'' (115.2).
The reassurance rests on two exact distinctions. Jesus did not say that his
kingdom is not \emph{in} this world but that it is not \emph{of} it, and not
that it is not \emph{here} but that it is not \emph{from hence}. ``For His
kingdom is here until the end of the world, having tares intermingled
therewith until the harvest,'' and yet ``it only sojourns as a stranger in the
world.'' The kingdom is ``those who believe in Him,'' who belonged to the
world while they belonged to its prince, and Augustine gives the Epistle's
verse for the change: ``For so did God rescue us from the power of darkness,
and translate us into the kingdom of the Son of His love'' (Col 1:13). ``Thou
sayest that I am a king'' neither denies a kingship nor confesses the one
Pilate imagines, ``Carnal thyself, thou sayest it carnally'' (115.3); and
``Every one that is of the truth heareth my voice'' commends ``the grace
whereby He calleth according to His own purpose'' (115.4). The kingdom gathers
its people by truth heard and grace given, not by the strength of servants.

St~John Chrysostom reads the same sentence of the kingdom's power rather than
its people. ``Not because He does not rule here, but because He has his empire
from above, and because it is not human, but far greater than this and more
splendid.'' The servants who do not fight are the proof: ``Here He shows the
weakness of kingship among us, that its strength lies in servants; but that
which is above is sufficient for itself, needing nothing'' (83.4). St~Thomas
Aquinas sets the two senses of \latin{regnum} side by side, the people ruled
with Augustine and the royal power with Chrysostom, cites the Alleluia's verse
for that power, and draws the conclusion that resolves the Mass's apparent
contradiction: \latin{non est hinc \dots\ est tamen hic, quia ubique est}, it
is not from here, and yet it is here, because it is everywhere (\work{Super
Io.}\ 18 lect.~6). A kingdom that took its origin from one place in the world
would be bounded by that place. One whose origin is above is bounded by none,
and can claim the whole earth without competing with any power in it. The
formulary keeps both of Aquinas's senses: the Collect asks for a people
gathered out of the scattered nations, the Alleluia sings a power that lasts.
Aquinas also glosses the Epistle's \latin{transtulit} as \latin{ut essemus
regnum Dei}, ``that we might be the kingdom of God,'' and sets the Gospel's
\latin{regnum meum non est de hoc mundo} beside it (\work{Super Col.}\ 1
lect.~3).

So the first and second readings differ at the Gospel without contradicting
each other. The first sees there a dominion received as man, won by the
passion that has begun. The second sees a power whose origin and means are not
from the world and which needs no servants to defend it; it still has the
Offertory give the nations as an inheritance. They meet at
Jn 18:37, where the second reading's own principal witnesses divide. Augustine
reads ``For this was I born'' of the birth from the Virgin (115.4), Chrysostom
of the eternal generation, and Aquinas gives both, the second \latin{sicut
Deus ex Deo, ita rex ex rege}, adding that ``for this came I into the world''
is the temporal coming even on Chrysostom's reading. The second reading rests
on the origin of the kingdom in verse~36 and on the hearers in verse~37, and
either birth serves it. For the first, Chrysostom's answer is its strongest
objection, and the first reading grants his premise and refers the receiving
to the humanity.
Augustine's own two loci speak of two births and do not conflict: he reads
Ps 2:7 of the eternal generation and Jn 18:37 of the birth from the Virgin,
so that the eternal birth grounds the decree of the psalm and the temporal
birth the testimony before Pilate.

\subsection{What the kings and nations owe}

A kingdom not from hence still asks a great deal of the world. The Gradual
sings that all kings shall adore him, the Offertory's psalm bids the kings
``serve ye the Lord with fear,'' the Collect asks that the nations be
subjected, \latin{subdantur}, and \work{Quas primas} calls it a grave error to
say that Christ has no authority in civil affairs, and draws from the feast a
public duty of rulers (nn.~17--19, 32). Read alone, the Gospel's denial might
seem to empty that subjection of all public consequence; read alone, the
subjection might seem to make Christ a rival of Caesar. The witnesses refuse
both, and they give the kings three distinct places.

The first is conversion. Augustine reads the Gradual's psalm of the Church
spread through the whole world, its river the Jordan from which the Gospel of
the kingdom went out (\work{Enarrationes in Psalmos} 71.11). When the kings of
Tharsis and Saba ``lead'' gifts, a verb for things that walk, the gifts ``seem
to me to signify men, whom into the fellowship of the Church of Christ the
authority of kings doth lead: although even persecuting kings have led gifts,
knowing not what they did, in sacrificing the holy Martyrs'' (71.13). And the
kings who adore belong to ``the people of men believing in Him,'' who do not
disdain to be poor and humbly confess their sins (71.14). St~Robert Bellarmine
reads the Offertory's psalm as an exhortation to ``the kings of this world
\dots\ to freely subject themselves to him, the true and supreme king of all
kings,'' serving as a son serves his father (\work{Explanatio}, Ps~2).
Cassiodorus states the paradox of the Gradual's verse: \latin{reges
adorabunt, qui putantur adorandi; gentes autem servire, quae terrenos dominos
habere noscuntur} (\work{Expositio Psalmorum} 71).

The second is office. In \work{Letter}~185, written to justify the imperial
laws against the Donatist schism, Augustine joins the Offertory's psalm to the
Gradual's verse himself. In the apostles' time no emperor served Christ by
laws, because the raging of the kings in Ps 2:1--2 was being fulfilled, not
yet the command of verses 10--11. ``For a man serves God in one way in that he
is man, in another way in that he is also king''; and once ``All kings shall
fall down before Him; all nations shall serve Him'' began to be fulfilled, no
sober man could tell kings that the Church was no concern of theirs
(185.19--20). Aquinas reports the distinction on Ps 2:10--11, \latin{secundum
Augustinum \dots\ rex servit Deo inquantum homo, in se juste vivendo, sed
inquantum rex, leges ferendo contra ea quae sunt contra Dei justitiam}, and
sees in the psalm the Church's history, the kings' laws first against Christ
and afterwards for him (\work{Postilla super Psalmos} 2, n.~9). The Introit's
verse names the measure of that office. Bellarmine gives its literal sense as
David's prayer for Solomon, who asks the grace of judging rightly, ``for kings
were created for the people, not the people for kings'' (\work{Explanatio},
Ps~71), while Augustine, Cassiodorus and Theodoret hold that the psalm fits
Solomon in no way and speaks of Christ alone, the King to whom the Father has
given all judgment (Jn 5:22).

The third is the judgment. Bellarmine finds the Gradual's adoration fulfilled
``to a certain extent'' in the conversion of ``a great many princes and
nations,'' ``to the letter'' at the day of judgment, when every knee bends, and
``always'' in ``the actual power that Christ has over all princes and all
nations, so that, with or without their knowledge, with or without their
consent, he may deal as he pleases with them'' (\work{Explanatio}, Ps~71).
Theodoret divides the verse the same way, some kings willingly in this life
and all after the resurrection, since ``now we see not as yet all things
subject to him'' (Heb 2:8). Schuster gives the state of things between the
two: the kingdom is ``complete and absolute de jure, but \dots\ opposed and
incomplete de facto'' (p.~202), and on the Gospel, ``Jesus does not come to
dethrone the sovereigns of the earth'' (p.~203).

The denial before Pilate, on this reading, is a denial of a worldly origin and
worldly means, not of authority over the world. The subjection is real,
rendered freely in conversion, rendered by rulers in their own office, and
completed at the judgment, over a dominion that depends on no one's consent.
None of the three moments makes Christ a competitor for Caesar's throne.
\work{Quas primas} holds the two halves in the same way: Christ during his
earthly life refrained from exercising such authority and does not now
interfere with those who possess earthly goods, and yet individuals, families
and states are under his rule, and rulers owe him public honour and obedience
(nn.~17--19, 32).

The Postcommunion presses the same question on the faithful. It is the one
text of the Mass that makes them soldiers, \latin{sub Christi Regis vexillis
militare}, a few minutes after the Gospel in which the King has said that if
his kingdom were of this world his servants would strive. Chrysostom's reasoning on Jn 18:36 shows what the figure cannot mean:
earthly kingship depends on servants who fight and Christ's does not, so their
warfare is not his defence. Hilary's reading of the second psalm shows what it
can mean: the kingdom of God is where sin is overcome and death destroyed, and
those who become co-heirs of the eternal King must themselves become kings
(\work{Tractatus super Psalmos} 2.42--43). The standard, in Schuster's note on
that prayer, is the Cross, ``for the kingdom of Christ is that of obedience,
humility and sacrifice'' (p.~204). The first reading hears in the same prayer
the members whom Bede makes sharers in the priesthood: \latin{Quia Rex regum et
Sacerdos coelestis se offerendo pro nobis nos suo corpori adunavit, nemo
sanctorum est qui spiritualiter sacerdotii officio careat} (\work{Explanatio
Apocalypsis} I, on 1:6), on the clause of John's greeting that the Introit
leaves out.

\subsection{A kingdom handed over that has no end}

The Preface follows Paul's order in 1~Cor 15: the Son anointed Priest and
King offers himself on the cross, all creatures are subjected to his rule,
and he hands over \latin{aeternum et universale regnum} to the Father. Paul's
own words stand behind the verb, ``when he shall have delivered up the kingdom
to God and the Father \dots\ For he must reign, until he hath put all his
enemies under his feet,'' and ``then the Son also himself shall be subject
unto him that put all things under him, that God may be all in all''
(15:24--25, 28). Against them the Alleluia sings a power that ``shall not be
taken away,'' the Communion a King who ``shall sit \dots\ for ever,'' and
Gabriel had told Mary, ``of his kingdom there shall be no end'' (Lk 1:33).

Some in St~Cyril of Jerusalem's day took the ``until'' to mean that the reign
stops. He warns his catechumens, ``shouldest thou ever hear any say that the
kingdom of Christ shall have an end, abhor the heresy; it is another head of
the dragon, lately sprung up,'' and sets against it Gabriel and Daniel's
everlasting power. Scripture's ``until,'' he argues, sets no term: death
reigned from Adam ``until'' Moses (Rom 5:14) and did not cease there. The Son's
subjection is no new obedience but ``a self-chosen accordance; for He is not a
servant, that He should be subjected by force, but a Son''
(\work{Catechetical Lectures} 15.27--32). Jerome had said the same from
Daniel: no man could hold a power \latin{quae nullo fine claudatur}, which no
end encloses (\work{In Dan.}\ II, on 7:14).

St~Augustine answers by asking what the kingdom is. ``When He shall have
delivered up the kingdom to God, even the Father'' is ``as much as to say,
When He shall have brought believers to the contemplation of God, even the
Father'' (\work{De Trinitate} I.8.16). ``For the Father will be manifested to
us, and that will suffice for us,'' and the heart is made ready for that sight
now: ``Blessed are the pure in heart, for they shall see God'' (I.8.17). The
handing over takes nothing from the Son, because what he hands over is not a
possession apart from himself: ``His faithful ones, whom He has redeemed with
His blood, are called His kingdom, for whom He now intercedes; but then,
making them to abide in Himself there, where He is equal to the Father, He
will no longer pray the Father for them'' (I.10.21). St~Hilary reaches the
same answer by the grammar of the verse. The Son delivers the kingdom ``not in
the sense that He resigns His power by the delivering, but that we, being
conformed to the glory of His body, shall form the Kingdom of God. It is not
said, He shall deliver up His Kingdom, but, He shall deliver up the Kingdom''
(\work{De Trinitate} XI.39). Augustine and Hilary both make the handing over
the vision of God, and both cite the same beatitude for the purity it
requires.

On the subjection that follows they speak differently. Augustine keeps the
form of the servant as what is subject: ``in so far as He is a Priest that has
taken the form of a servant for us,'' he ``shall be put under Him who has put
all things under Him'' (I.10.20). Hilary writes that the end of the subjection
is ``that God may be all in all, that no trace of the nature of His earthly
body may remain in Him. Although before this time the two were combined within
Him, He must now become God only; not, however, by casting off the body, but
by translating it through subjection'' (XI.40). He insists that the body is
translated and not cast off, while his ``God only'' goes further than
Augustine's servant who remains subject. The Preface uses neither
formulation; it says only that the kingdom is handed over and that it is
eternal.

The third reading draws the rest of the Mass toward that end. Hilary says, on
the Offertory's psalm, that the nations given to the Son are the eternal
inheritance of the firstborn from the dead (\work{Tractatus super Psalmos}
2.31), and, in the \work{De Trinitate}, that the kingdom the Son delivers is
``us who have been made the Kingdom'' (XI.39); the reading puts the two
together, so that what is given to the Son is what the Son gives to the
Father. The Epistle's firstborn from the dead reconciles things ``on earth''
and ``in heaven,'' which Chrysostom explains by Christ's bringing man up to
the angels and joins to the summing up of all things in Eph 1:10, the
Collect's \latin{instaurare} (\work{Homilies on Colossians} 3). The Secret shows
the Priest still offering and still giving unity and peace now; Augustine says
that when the kingdom is delivered up he will no longer intercede for us as
Priest (I.10.20). And the Gospel's \latin{non est hinc} names a kingdom whose
origin is the Father's and whose end is the Father's too.

\subsection{When the King sits}

The Communion's psalm carries, in the verse the antiphon cuts, the flood that
the Lord ``maketh \dots\ to dwell,'' and the witnesses' readings of the flood
place the session for ever at different moments. St~Thomas gives first the
plain sense, God's judgments manifest from the beginning, as in the deluge,
and then three mystical ones: the saints secure in the Church as in Noe's ark,
the world and its carnal men brought into his temple by conversion, \latin{et
sedebit rex in aeternum}, and the baptismal waters (\work{Postilla super
Psalmos} 28). St~Augustine reads the flood as this world: ``The Lord therefore
first inhabiteth the deluge of this world in His Saints, kept safely in the
Church, as in the ark. `And the Lord shall sit a King for ever.' And afterward
He will sit reigning in them for ever'' (\work{Enarrationes in Psalmos}
28.10). The blessing in peace is ``affording them peace in Himself,'' ``for
peace in this world He hath not promised them'' (28.11). St~Basil the Great
reads the flood as the grace of baptism, after which God, settled in the soul
that shines from that deluge, makes it, as it were, a throne for himself
(\work{Homilia in Psalmum} 28.8). Blessed Ildefonso Schuster reads the
antiphon of the communicant: ``Jesus is enthroned in the faithful heart''
(p.~204). St~Robert Bellarmine, on his second reading, places the session
after the judgment: ``all resistance being broken down, the whole power of
demons, bad men, and all power in general being swept away, `the Lord shall
sit King forever' '' (\work{Explanatio}, Ps~28).

The session therefore has a beginning now, in the soul cleansed by baptism
(Basil), in the communicant (Schuster), in the saints amid the flood of this
world (Augustine), and a completion after the judgment. Augustine's session
runs from the one into the other; Bellarmine's begins when the judgment is
done. The Postcommunion stands between them: it is said by those who have
just received \latin{immortalitatis alimoniam}, and it asks to reign with him
\latin{in caelesti sede}. Each reading hears the Communion in its own key. For
the first it is the session of the exalted King who received the kingdom; for
the second its peace is the peace in himself that Augustine distinguishes
from the world's; for the third its ``for ever'' is the reign that the
handing over does not end.

\subsection{Three answers and their senses}

The three readings share more than they divide. All identify the kingdom with
the redeemed people, the Church, and none lets it become a territory. The
cross has a place in each, and a different one. In the first it is the means
by which the kingdom is received, the asking of the second psalm, offered
again in the Secret's victim. In the second it is the standard under which the
faithful fight, in Schuster's note, and their warfare is against sin. In the
third it is the priestly offering that the Preface sets before the handing
over. Augustine contributes to all three readings, and so does Hilary.

Their four senses diverge where their questions do. In the literal sense, the
first reading hears royal psalms that promise the nations to a Son who asks,
and a Lamb acclaimed because he was slain; it keeps Bellarmine's literal
Solomon beside the Christ alone of Augustine, Cassiodorus and Theodoret. The
second hears a prisoner who denies his kingdom a worldly origin and worldly
defenders, and psalms that warn the kings to serve. The third hears Daniel's
everlasting kingdom, Paul's handing over, and the psalm's flood and session.
Allegorically, the first finds the Church as the inheritance the Son receives
as man; the second, the Church here and not from here, spread from the Jordan;
the third, the Church as the kingdom the Son hands over. Morally, the first
asks the faithful to ask and to share the King's self-offering; the second, to
hear the truth and serve him in their own office, against sin and not with
swords; the third, the purity of heart that the vision requires. And
anagogically the first sees the firstborn from the dead go before his members
into the reign; the second, the mixed kingdom of the tares made manifest at
the harvest, when every knee bends; the third, the vision face to face, God
all in all, and a King not diminished by giving his people to the Father.

Each answer would mislead on its own. The first, alone, would mislead if it
let the kingship begin at Easter. The second, alone, could leave the Preface's
handing over and the Communion's eternal session as ornaments of a kingdom
wholly present in history. The third, alone, could pass over the Gospel's
present, mixed kingdom and the service that kings owe now. Together they say
what the Preface says in one sentence: the anointed Son reigns by right,
reigns as the crucified Priest, and reigns towards the Father, in whom his
kingdom of truth and life, holiness and grace, justice, love and peace has no
end.
